\documentclass[11pt,a4paper]{article}
\usepackage{amsmath,amssymb,amsthm}
\usepackage[colorlinks=true,linkcolor=blue,citecolor=blue,urlcolor=blue]{hyperref}
\usepackage[margin=1in]{geometry}
\usepackage{xcolor}
\usepackage{booktabs}
\usepackage{array}
\usepackage{float}
\usepackage{mathtools}

\newcommand{\Nanti}{N_{\mathrm{anti}}}
\newcommand{\Sp}{\mathrm{Sp}}
\newcommand{\rad}{\mathrm{rad}}
\newcommand{\rank}{\mathrm{rank}}
\newcommand{\F}{\mathbb{F}}
\newcommand{\Z}{\mathbb{Z}}

\theoremstyle{plain}
\newtheorem{theorem}{Theorem}[section]
\newtheorem{lemma}[theorem]{Lemma}
\newtheorem{corollary}[theorem]{Corollary}
\newtheorem{proposition}[theorem]{Proposition}
\theoremstyle{definition}
\newtheorem{definition}[theorem]{Definition}
\theoremstyle{remark}
\newtheorem*{remark}{Remark}

\title{\textbf{The Arity--Resonance Ceiling:\\
  Why Pauli Contextuality Stops at $H^3$}\\[4pt]
\large (Paper XXII)}
\author{Zhou-Li Chen\\[2pt]\small co-nlang Research}
\date{June 2026}

\begin{document}
\maketitle

\begin{abstract}
Papers XVIII--XXI located the Mermin obstruction in $H^3$ of the
$5$-Lagrangian nerve and proved it vanishes universally if and only if $n=4$.
A natural worry remained: if the obstruction sits in $H^3$, why not $H^4$, $H^5$,
$\ldots$?  This paper answers it with an \emph{arity--resonance principle}.
A relative-position datum of arity~$a$ is intrinsically an $(a-1)$-cochain on the
index nerve; on a proper $K_N$ (nerve $\partial\Delta^{N-1}\cong S^{N-2}$) it lands
in the top cohomology iff $N=a+1$.  The Pauli data supplies exactly two natural
cochains---the Maslov bit $\mu$ (arity~$3$) and the anticommutation cochain
$\mathbf a$ (arity~$4$, with $\mathbf a$ a coboundary $\delta\mu$ at $n=4$)---so the
tower has two live rungs, $K_4/H^2$ and $K_5/H^3$, and a \emph{lid}.
We prove the truncation: on $K_6$ the anticommutation cochain is sub-top, its
degree-$4$ assembly is the exact coboundary $\delta\mathbf a$, and the
$H^4$ class vanishes for every proper $K_6$ (which we show exist).  We package the
whole tower into one graph invariant: the family-$A$ cohomological ceiling equals
$\min(\omega(G)-2,\,3)$, where $\omega(G)$ is the clique number of the incidence
graph; the pentagram ($\omega=5$) is the unique resonance.  The $n$-dependence of
each rung is governed by the single rank-parity lemma of Paper~XX.
Finally, the Mermin--Peres square---six Lagrangians in $\Sp(4,\F_2)$ with bipartite
($K_{3,3}$) incidence---shows the ceiling is configuration-specific: it carries no
family-$A$ obstruction and lives instead in the central-extension class, exhibiting
a second contextuality family.
\end{abstract}

\section{Introduction}

Quantum contextuality is a local-to-global obstruction: the data of individual
measurement contexts cannot always be glued into one consistent global assignment,
and cohomology measures the failure.  Paper~XX~\cite{PaperXX} showed that
$\Nanti\bmod2$ is the evaluation of a degree-$3$ class on the nerve of a proper
$K_5$, and Paper~XXI~\cite{PaperXXI} proved it vanishes universally exactly at
$n=4$.  Both papers live in degree~$3$.
In an ordinary obstruction ladder one expects the story to continue:
$H^3\to H^4\to H^5\to\cdots$.  It does not, and the reason is the point of this
paper.

The degree at which the obstruction lives is not free: it is pinned by the
\emph{arity} of the algebraic datum that builds it.  Anticommutation
$\omega(v_{ij},v_{kl})$ involves four indices, so it is a $3$-cochain---no matter how
many contexts there are.  On the nerve of a proper $K_5$, which is a $3$-sphere,
degree~$3$ is exactly the top, and the obstruction is genuine.  On a proper $K_6$,
whose nerve is a $4$-sphere, the same datum is one degree below the top and is
forced to be a coboundary.  The pentagram is the unique configuration size at which
arity-$4$ data resonates with the top of the nerve.

\paragraph{Results.}
\begin{enumerate}
\item \textbf{The arity--resonance principle} (\S\ref{sec:resonance}): an arity-$a$
  datum is an $(a-1)$-cochain; it hits the top cohomology of a proper $K_N$
  ($\cong S^{N-2}$) iff $N=a+1$.
\item \textbf{Two rungs} (\S\ref{sec:rungs}): the Maslov bit ($a=3$) resonates at
  $K_4/H^2$; the anticommutation cochain ($a=4$) resonates at $K_5/H^3$.  The
  $n$-dependence of each follows from the rank-parity lemma of Paper~XX:
  $H^2$ opens at odd $n\ge5$, $H^3$ at all $n\ge5$.
\item \textbf{The truncation theorem} (\S\ref{sec:trunc}): proper $K_6$ exist, but
  the degree-$4$ assembly of the anticommutation cochain is the exact coboundary
  $\delta\mathbf a$, so the $H^4$ class vanishes universally.  The obstruction
  truncates at $H^3$.
\item \textbf{The clique criterion} (\S\ref{sec:clique}): the family-$A$ ceiling is
  $\min(\omega(G)-2,3)$, with $\omega(G)$ the clique number of the incidence graph.
\item \textbf{Two families} (\S\ref{sec:square}): the Mermin--Peres square is
  bipartite ($\omega=2$); it carries no family-$A$ obstruction and realizes the
  central-extension class, a distinct family.
\end{enumerate}

\paragraph{Conventions.}
$V=\F_2^{2n}$ with the standard symplectic form $\omega$.  A \emph{proper $K_N$} is
an $N$-tuple of Lagrangians $(L_0,\dots,L_{N-1})$ with $\dim(L_i\cap L_j)=1$ for all
$i\ne j$ and all $\binom N2$ rays $v_{ij}:=L_i\cap L_j$ distinct.

\section{Setup}\label{sec:setup}

\subsection{The nerve and its cochains}

The index set $\{0,\dots,N-1\}$ spans an abstract simplex $\Delta^{N-1}$.  Because a
proper $K_N$ has no $N$-fold common ray (the configuration is ``hollow''), the
relevant complex is the boundary $\partial\Delta^{N-1}\cong S^{N-2}$, with
\[
  H^k(\partial\Delta^{N-1};\F_2)=\begin{cases}\F_2 & k\in\{0,\,N-2\}\\ 0 &\text{otherwise.}\end{cases}
\]
Two relative-position cochains are canonically attached to the configuration.

\begin{definition}[Maslov $2$-cochain]
For a triple $\{a,b,c\}$, let $\mu(a,b,c)\in\F_2$ be the bit recording whether
$\omega\not\equiv0$ on the triple data (Paper~XX~\cite{PaperXX}); it is
$S_3$-symmetric and $\Sp$-invariant.
\end{definition}

\begin{definition}[Anticommutation $3$-cochain]
For a $4$-index set $T=\{i,j,k,l\}$ let $\mathbf a_T\in\F_2$ be the parity of the
number of anticommuting cross-context pairs among the three pairings
$\{v_{ij},v_{kl}\},\{v_{ik},v_{jl}\},\{v_{il},v_{jk}\}$.  This defines a $3$-cochain
$\mathbf a$ on the nerve; on a proper $K_5$, $\Nanti\bmod2=\langle\mathbf a,[S^3]\rangle$.
\end{definition}

\subsection{The incidence graph}

\begin{definition}[Incidence graph]
The incidence graph $G$ of a Lagrangian configuration has the contexts as vertices
and an edge $\{i,j\}$ whenever $\dim(L_i\cap L_j)=1$.  Write $\omega(G)$ for its
clique number.  A proper $K_N$ has complete incidence, $\omega=N$.
\end{definition}

\section{The arity--resonance principle}\label{sec:resonance}

The degree of a contextuality obstruction is not an abstract label: it counts how
many contexts must be observed \emph{together} before the obstruction is visible at
all.  A datum that depends on $a$ contexts can only be recorded on the
$(a-1)$-faces of the index nerve, so its natural degree is $a-1$.  The obstruction
becomes genuine precisely when this degree reaches the top of the nerve --- a
resonance between the arity of what can be jointly observed and the size of the
configuration.

\begin{proposition}[Arity is degree]\label{prop:arity}
A relative-position invariant that depends on $a$ contexts is a cochain of degree
$a-1$ on the index nerve.  On a proper $K_N$ ($\cong S^{N-2}$) its class lies in
$H^{a-1}(S^{N-2};\F_2)$, which is nonzero (and equals the top) iff $a-1=N-2$, i.e.
\[
  N=a+1 .
\]
\end{proposition}

\begin{proof}
An $a$-ary datum is a function of an $a$-element index subset, i.e. of an
$(a-1)$-face; it is therefore an $(a-1)$-cochain.  By the cohomology of
$S^{N-2}$ above, the ambient cohomology \emph{group} $H^{a-1}(S^{N-2};\F_2)$ is
nonzero only for $a-1\in\{0,N-2\}$; the top (and only positive-degree) value is
$a-1=N-2$.
\end{proof}

This proposition is about \emph{where the obstruction can live}, i.e. the
nonvanishing of the ambient group; that the specific class built from the Pauli
data is in fact nontrivial at the resonant $N$ is the content of
Papers~XX--XXI~\cite{PaperXX,PaperXXI}, not of this proposition.

Thus each arity has a \emph{resonant configuration} $K_{a+1}$.  Below it
($N<a+1$) the datum is over-determined; above it ($N>a+1$) it is sub-top and
cohomologically trivial.  The Pauli data realizes two arities.

\begin{remark}
The symplectic form is bilinear (arity $2$); the Maslov triple raises this to
arity $3$ and the anticommutation pairing of rays to arity $4$, each extra index
supplied by the \emph{configuration} rather than by the form.  That this is the top
is motivated, though not proved, by two facts: the canonical degree-$4$ assembly of
the anticommutation data is the \emph{exact} coboundary $\delta\mathbf a$
(Theorem~\ref{thm:trunc}), carrying no new class; and the natural quadruple bit
$q_4$ of Paper~XIX~\cite{PaperXIX} is itself arity-$4$ (a $3$-cochain) and
empirically saturated.  A genuinely new arity-$5$ invariant would have to be a datum
\emph{not} assembled from the bilinear form; we conjecture none exists, so the
tower's top rung is $a=4$ (see also the rigor table, where this is the one
conjectural entry).
\end{remark}

\section{The two rungs}\label{sec:rungs}

The Pauli data supplies two cochains, and they are one phenomenon seen at two
configuration sizes: the Maslov bit resonates with the $4$-context nerve $S^2$, the
anticommutation cochain with the $5$-context nerve $S^3$.  Their $n$-dependence is
governed by a single lemma --- the rank-parity lemma of Paper~XX --- which we
recall as each rung is reached.

\subsection{$K_5/H^3$: the anticommutation rung}

On a proper $K_5$ the nerve is $S^3$, whose top cells are the five tetrahedra;
$\mathbf a$ is automatically a cocycle (no $4$-cells), and
$\langle\mathbf a,[S^3]\rangle=\Nanti\bmod2$ is the $H^3$ class
(Papers~XX--XXI).  By the master theorem of Paper~XXI~\cite{PaperXXI} it vanishes
universally iff $n=4$; at $n=3$ the all-Fano subclass gives $\Nanti=15$, and at
$n\ge5$ it is generically nontrivial.

\subsection{$K_4/H^2$: the Maslov rung}

On a proper $K_4$ the nerve is $\partial\Delta^3\cong S^2$, whose top cells are the
four triangles; $\mu$ is a $2$-cocycle and its class is
$\langle\mu,[S^2]\rangle=\sum_{\{a,b,c\}}\mu(a,b,c)\bmod2$.

\begin{proposition}[The Maslov rung]\label{prop:k4}
For a proper $K_4$ in $\Sp(2n,\F_2)$ the Maslov $H^2$ class
$\langle\mu,[S^2]\rangle$ \emph{vanishes} for $n=3$ and for all even $n$.  For odd
$n\ge5$ it is nontrivial (both values occur), computationally.
\end{proposition}

\begin{proof}
By the rank-parity lemma of Paper~XX~\cite{PaperXX} ($\rank B=n-3$ for the symmetric
form $B$ governing $\mu$): $\mu\equiv0$ at $n=3$ and $\mu\equiv1$ for even $n$.  In
both cases the four summands are uniformly $0$ or uniformly $1$, so
$\langle\mu,[S^2]\rangle=\sum_{\text{4 triangles}}\mu=0$ --- this part is rigorous.
For odd $n\ge5$ the lemma only says $\mu$ varies \emph{per triple}; that the sum
$\langle\mu,[S^2]\rangle$ then attains \emph{both} values across proper $K_4$ is
verified computationally (\texttt{resonance\_tower.py}~\cite{PaperXXIIsuppl}: at
$n=5$, $211$ vs $189$ of $400$ sampled $K_4$).  Rank-parity alone does not exclude
a uniform value at some odd $n$, so the opening is recorded as a numerical, not
algebraic, claim.
\end{proof}

\begin{corollary}[The tower's $n$-dependence]\label{cor:ndep}
\begin{center}
\renewcommand{\arraystretch}{1.2}
\begin{tabular}{lll}
\toprule
rung & class & opens at \\
\midrule
$K_4/H^2$ (Maslov) & $\langle\mu,[S^2]\rangle$ & odd $n\ge5$ \\
$K_5/H^3$ (anticomm.) & $\Nanti\bmod2$ & all $n\ge5$ \;(and $n=3$ all-Fano)\\
\bottomrule
\end{tabular}
\end{center}
A single lemma (rank parity) governs both: $H^2$ is a pure-$\mu$ phenomenon and opens
only where $\mu$ varies; $H^3$, carried by $\mathbf a$, has the additional
even-$n$ channel of Paper~XXI and opens at every $n\ge5$.
\end{corollary}

\section{The truncation theorem}\label{sec:trunc}

We now show the obstruction stops at $H^3$.  The mechanism is the arity ceiling:
anticommutation is a $3$-cochain, top-dimensional on the pentagram's $S^3$ but
\emph{sub}-top on the next configuration's $S^4$, where it is forced to be exact.
Concretely, the only degree-$4$ quantity one can assemble from it turns out to be a
coboundary --- so it carries no class, and proper $K_6$ (which exist) have no
$H^4$ obstruction.

\begin{proposition}[Proper $K_6$ exist]\label{prop:k6exist}
For every $n\ge4$ there is a proper $K_6$ in $\Sp(2n,\F_2)$.
\end{proposition}

\begin{proof}
In the chart of Lagrangians transverse to a fixed $L_\infty$, $L_S=\{(x,Sx)\}$ for
symmetric $S$, one has $L_S\cap L_{S'}=\ker(S+S')$; six symmetric $n\times n$
matrices with pairwise sums of rank $n-1$ and distinct kernels give a proper $K_6$.
Such sextuples are exhibited computationally for $n=4,5,6$
(\texttt{k6\_truncation.py}~\cite{PaperXXIIsuppl}); for larger $n$ they persist by
the spread-stabilisation lemma of Paper~XXI~\cite{PaperXXI} (orthogonal direct sum
with a trivial $\Sp(2m)$ block, appending pairwise-transverse Lagrangians).
\end{proof}

\begin{remark}[An explicit sextuple]
For $n=4$ the six symmetric matrices below have all pairwise sums $S_i+S_j$ of
rank~$3$ (so each $L_{S_i}\cap L_{S_j}$ is a single ray), with the $15$ rays
distinct --- hence a proper $K_6$ in $\Sp(8,\F_2)$:
\[
S_0=\begin{psmallmatrix}1&1&1&1\\1&0&1&0\\1&1&1&1\\1&0&1&0\end{psmallmatrix}\!,\quad
S_1=\begin{psmallmatrix}0&0&1&1\\0&1&0&1\\1&0&0&1\\1&1&1&1\end{psmallmatrix}\!,\quad
S_2=\begin{psmallmatrix}0&1&0&0\\1&1&0&1\\0&0&0&1\\0&1&1&0\end{psmallmatrix},
\]
\[
S_3=\begin{psmallmatrix}1&1&1&1\\1&1&1&0\\1&1&0&0\\1&0&0&0\end{psmallmatrix}\!,\quad
S_4=\begin{psmallmatrix}0&0&1&1\\0&0&1&0\\1&1&0&1\\1&0&1&0\end{psmallmatrix}\!,\quad
S_5=\begin{psmallmatrix}0&0&1&0\\0&1&0&0\\1&0&1&1\\0&0&1&1\end{psmallmatrix}.
\]
\end{remark}

\begin{theorem}[Truncation at $H^3$]\label{thm:trunc}
On a proper $K_6$ the nerve is $\partial\Delta^5\cong S^4$.  Define the degree-$4$
cochain $\mathbf c$ by $\mathbf c_m=\Nanti(\text{face }m)\bmod2$, where face~$m$ is
the $K_5$ on the five indices $\ne m$.  Then $\mathbf c=\delta\mathbf a$, the
coboundary of the anticommutation $3$-cochain; hence $[\mathbf c]=0$ in
$H^4(S^4;\F_2)$, and the anticommutation data carries no $H^4$ obstruction.
\end{theorem}

\begin{proof}
Face~$m$ is a $4$-simplex of $\partial\Delta^5$; its five $3$-faces are the
$4$-index subsets $T\subset\{0,\dots,5\}\setminus\{m\}$.  Each cross-context pair
$\{v_{ij},v_{kl}\}$ uses exactly four indices $\{i,j,k,l\}$, so the $15$
cross-context pairs of the $K_5$ face~$m$ are \emph{partitioned} by its five
$4$-index subfacets $T$ (each pair lies in the unique $T$ equal to its index set);
hence by definition $\Nanti(\text{face }m)=\sum_{T\subset\text{face }m}\mathbf a_T$, so
\[
  \mathbf c_m=\sum_{T:\,T\subset\partial(\text{face }m)}\mathbf a_T
            =\langle\mathbf a,\partial(\text{face }m)\rangle
            =(\delta\mathbf a)_m .
\]
Thus $\mathbf c=\delta\mathbf a$ is a coboundary and $[\mathbf c]=0$.  Equivalently,
$\langle\mathbf c,[S^4]\rangle=\langle\mathbf a,\partial[S^4]\rangle=0$ since $[S^4]$
is a cycle---each cross-context pair lies in exactly two of the six faces.
\end{proof}

\begin{remark}[Why $K_5$ is the ceiling]
On $S^3$ ($K_5$) the anticommutation $3$-cochain is top-dimensional, hence a cocycle,
and represents a genuine $H^3$ class.  On $S^4$ ($K_6$) it is one degree below the
top: it is no longer forced to be a cocycle, but its class lies in
$H^3(S^4;\F_2)=0$, and its only degree-$4$ shadow, $\delta\mathbf a$, is exact.
The same arity-$4$ datum is resonant at $K_5$ and silent at $K_6$.  This is the
truncation: the obstruction does not climb, because the data does not.
\end{remark}

\begin{remark}[Live $H^3$, dead $H^4$]
The vanishing of $[\mathbf c]$ is not a scarcity of structure.  Over $200$ proper
$K_6$ at $n=6$, the six face-classes $(\mathbf c_m)\in\F_2^6$ realize all $32$
even-weight patterns---the full kernel of $\Sigma\colon\F_2^6\to\F_2$, the $H^4$
pairing---so the sub-pentagram $H^3$ classes are maximally free, constrained only by
$\sum_m\mathbf c_m=0$ (\texttt{k6\_truncation.py}~\cite{PaperXXIIsuppl}).
\end{remark}

\section{The clique criterion}\label{sec:clique}

The resonance tower compresses into a single graph invariant.

\begin{theorem}[Incidence-clique ceiling]\label{thm:clique}
Let $G$ be the incidence graph of a Lagrangian configuration.  The family-$A$
(anticommutation/Maslov) nerve obstruction is supported on the clique complex of
$G$; a degree-$d$ class requires a hollow $(d+2)$-clique.  Hence the family-$A$
ceiling is
\[
  \min\bigl(\omega(G)-2,\ a_{\max}-1\bigr)=\min(\omega(G)-2,\ 3),
\]
nontrivial only for $\omega(G)\ge3$, and saturating the arity-$4$ ceiling at
$\omega(G)=5$ (the pentagram).
\end{theorem}

\begin{proof}
A nerve cochain of degree $d$ is supported on $d$-faces, which in the clique complex
are $(d+1)$-cliques of $G$; a nonzero hollow class needs the boundary of a
$(d+2)$-clique (an $S^{d}$).  Thus the maximal achievable degree is $\omega(G)-2$.
It is further capped by the arity of the available data: the top Pauli cochain has
arity $a_{\max}=4$, degree $3$ (\S\ref{sec:resonance}).
\end{proof}

\begin{center}
\renewcommand{\arraystretch}{1.15}
\begin{tabular}{lccl}
\toprule
configuration & incidence & $\omega(G)$ & family-$A$ ceiling \\
\midrule
Mermin square & $K_{3,3}$ (triangle-free) & $2$ & none (family $B$) \\
proper $K_4$ & complete & $4$ & $H^2$ \\
proper $K_5$ (pentagram) & complete & $5$ & $H^3$ (resonance) \\
proper $K_6$ & complete & $6$ & $H^3$ (capped; $H^4$ exact) \\
\bottomrule
\end{tabular}
\end{center}
\noindent Verified in \texttt{clique\_criterion.py}~\cite{PaperXXIIsuppl}.

\section{A second family: the Mermin--Peres square}\label{sec:square}

The ceiling is configuration-specific.  The Mermin--Peres magic square is, in our
language, six Lagrangians in $\Sp(4,\F_2)$: three ``rows'' and three ``columns'',
\[
\begin{array}{ccc}
XI & IX & XX\\ IZ & ZI & ZZ\\ XZ & ZX & YY
\end{array}
\]
each row and column an isotropic plane (a context of three commuting observables).

\begin{proposition}[The square is bipartite, family $B$]\label{prop:square}
The incidence graph of the Mermin--Peres square is the complete bipartite graph
$K_{3,3}$: rows are pairwise transverse, columns are pairwise transverse, and each
row meets each column in one ray (the nine observables).  Its clique number is
$\omega=2$, so by Theorem~\ref{thm:clique} it carries \emph{no} family-$A$
obstruction; equivalently every triple of contexts contains a transverse pair, so
$\mu$ and $\mathbf a$ are undefined.  Its contextuality is instead the
$\pm I$ sign mismatch (one column product is $-I$): the central-extension /
quadratic-refinement class of Papers~III--VIII.
\end{proposition}

\begin{proof}
Direct computation of the six planes and their pairwise intersections gives the
$K_{3,3}$ incidence and nine distinct rays; all $20$ context-triples
contain a transverse pair.  Tracking Hermitian-Pauli phases, the six context
products are $+I$ except the $YY$ row, whose product is $-I$; the global product
$\prod_{\text{contexts}}(\pm I)=-I$ is the obstruction
(\texttt{mermin\_square.py}~\cite{PaperXXIIsuppl}).
\end{proof}

Thus contextuality configurations split, at minimum, into two cohomological
families: the all-pairwise family ($\omega\ge3$), whose nerve/anticommutation
obstruction climbs the resonance tower to its ceiling $H^3$ at the pentagram; and
the bipartite/triangle-free family ($\omega=2$), whose obstruction is the
central-extension sign class.  The pentagram's $H^3$ is the ceiling of the former,
not a universal one.

\section{Outlook}

Two natural questions are deferred.  First, whether the bipartite family has its own
resonance---a larger $K_{m,n}$ whose central-extension class reaches degree~$3$.
Second, whether the two families are pages of a single spectral sequence: the
Lyndon--Hochschild--Serre sequence of the Heisenberg extension
$1\to\Z_2\to P_n\to V\to1$ has the symplectic form $\omega\in H^2(V)$
(the family-$B$ class) as its transgression, and $\mathrm{Sq}^1\omega\in H^3(V)$ as
the natural candidate for the family-$A$ source~\cite{PaperIV}.  A direct
Maslov-side realization fails---the cup-$1$ square $\mu\cup_1\mu$ does not reproduce
$\Nanti\bmod2$, because $\mu$ is not a cocycle off the resonance
(\texttt{geometric\_route.py}~\cite{PaperXXIIsuppl})---so the unification requires
the homotopy-coherent comparison map $\partial\Delta^4\to BV$ whose coherence
obstruction is precisely $\mathbf a$.  This is the subject of future work.

\section*{Computational reproducibility}

All claims are checked by pure-Python scripts in the supplementary
archive~\cite{PaperXXIIsuppl}: \texttt{resonance\_tower.py} (the $K_4/H^2$ rung),
\texttt{k6\_truncation.py} (proper $K_6$, $\mathbf c=\delta\mathbf a$, the
even-weight saturation), \texttt{clique\_criterion.py} (the ceiling invariant),
\texttt{mermin\_square.py} (the bipartite family), and
\texttt{geometric\_route.py} (the negative cup-$1$ test).

\section*{Rigor self-assessment}

\begin{center}
\renewcommand{\arraystretch}{1.1}
\begin{tabular}{p{6.6cm}ll}
\toprule
Claim & Status & Source \\
\midrule
Arity $a$ $=$ degree $a-1$; resonance $N=a+1$ & Rigorous & Prop.~\ref{prop:arity} \\
$K_4/H^2$ rung: vanishing ($n=3$, even $n$) & Rigorous & Prop.~\ref{prop:k4} (XX rank-parity) \\
$K_4/H^2$ rung: opening (odd $n\ge5$) & Computational & \texttt{resonance\_tower.py} \\
$K_5/H^3$ rung & Rigorous & XX, XXI \\
Proper $K_6$ exist & Rigorous & Prop.~\ref{prop:k6exist} \\
$\mathbf c=\delta\mathbf a$, $[\mathbf c]=0$ (truncation) & Rigorous & Thm.~\ref{thm:trunc} \\
Clique ceiling $\min(\omega-2,3)$ & Rigorous & Thm.~\ref{thm:clique} \\
Mermin square bipartite, family $B$ & Rigorous & Prop.~\ref{prop:square} \\
No arity-$5$ invariant (the lid) & Conjectural & $q_4$ saturated, XIX \\
Even-weight saturation at $K_6$ & Empirical & \texttt{k6\_truncation.py} \\
\bottomrule
\end{tabular}
\end{center}

\begin{thebibliography}{99}

\bibitem[PaperIII]{PaperIII}
Z.-L. Chen, ``Kochen--Specker Contextuality as Central Extension:
The Peres--Mermin Square and the Pauli Group,''
\textit{Zenodo, DOI: 10.5281/zenodo.20073127} (2026).

\bibitem[PaperIV]{PaperIV}
Z.-L. Chen, ``The Cohomological Obstruction Ladder:
Lyndon--Hochschild--Serre Transgression and the $H^3$ Frontier,''
\textit{Zenodo, DOI: 10.5281/zenodo.20073184} (2026).

\bibitem[PaperVIII]{PaperVIII}
Z.-L. Chen, ``The $\Phi$ Functor and the $n$-Qubit Obstruction Ladder:
Explicit Transgression and Twisted Penrose Transform,''
\textit{Zenodo, DOI: 10.5281/zenodo.20454120} (2026).

\bibitem[PaperXV]{PaperXV}
Z.-L. Chen, ``The Weil Representation and the $S_3$ Lifting Classification
for Mermin Pentagrams,''
\textit{Zenodo, DOI: 10.5281/zenodo.20519654} (2026).

\bibitem[PaperXVIII]{PaperXVIII}
Z.-L. Chen, ``Mermin Pentagrams in $\Sp(8,\F_2)$:
Universal $N_{\mathrm{anti}}=10$ and the Gram Rank Selection Principle,''
\textit{Zenodo, DOI: 10.5281/zenodo.20579239} (2026).

\bibitem[PaperXIX]{PaperXIX}
Z.-L. Chen, ``Mermin Pentagrams in $\Sp(2n,\F_2)$, $n\ge5$:
Structure, Obstruction, and the Modulus Phenomenon,''
\textit{Zenodo, DOI: 10.5281/zenodo.20590195} (2026).

\bibitem[PaperXX]{PaperXX}
Z.-L. Chen, ``$H^3$ Opens at $n\ge5$:
The Cohomological Origin of the Mermin Modulus,''
\textit{Zenodo, DOI: 10.5281/zenodo.20608302} (2026).

\bibitem[PaperXXI]{PaperXXI}
Z.-L. Chen, ``$H^3$ Rigidity is Unique at $n=4$:
Master Theorem and the Even/Odd Obstruction Dichotomy,''
\textit{Zenodo, DOI: 10.5281/zenodo.20685669} (2026).

\bibitem[PaperXXIIsuppl]{PaperXXIIsuppl}
Z.-L. Chen, ``Supplementary computational scripts for Paper~XXII,''
\textit{Zenodo, DOI: 10.5281/zenodo.20685721} (2026).

\end{thebibliography}

\end{document}
